\chapter[Introdução]{Introdução}
O presente capítulo tem como objetivo apresentar as considerações iniciais sobre o Trabalho de Conclusão de Curso, que visa o desenvolvimento de uma aplicação web eficiente e centralizada para o gerenciamento de estágios na Universidade de Brasília (UnB). Inicialmente, será realizada uma Contextualização,
abordando o domínio de interesse deste trabalho, focado na gestão de estágios, e destacando a importância de um sistema centralizado que integre todas as etapas do processo de estágio, desde a inscrição até a conclusão.

Em seguida, será apresentada a Justificativa para a realização deste trabalho, esclarecendo as contribuições que a aplicação traz ao domínio da gestão acadêmica na UnB, especialmente no que tange à eficiência operacional,
acessibilidade e transparência dos dados relacionados aos estágios dos alunos.A proposta de desenvolvimento será orientada por princípios de desenvolvimento ágil, visando a construção de uma plataforma intuitiva e de fácil manutenção.


\section{Justificativa}

Atualmente, o gerenciamento dos estágios na FCTE é realizado de forma fragmentada e majoritariamente manual, envolvendo o uso de múltiplos sistemas, planilhas e trocas de e-mails entre alunos, docentes e coordenadores de curso. Esse processo descentralizado acarreta uma série de problemas, como a redundância de dados, inconsistência nas informações, perda de prazos, dificuldades na fiscalização e acompanhamento das atividades de estágio, além da sobrecarga administrativa para os setores responsáveis.

A ausência de uma ferramenta centralizada compromete a eficiência e a rastreabilidade das etapas que envolvem a execução dos estágios obrigatórios e não obrigatórios. Além disso, os alunos enfrentam incertezas e atrasos devido à falta de padronização nos procedimentos e à dificuldade de acesso às informações atualizadas sobre o andamento de seus processos.

Diante desse cenário, a implementação de um Sistema de Gerenciamento de Estágios surge como uma solução tecnológica capaz de integrar as diferentes etapas e agentes envolvidos no processo, promovendo maior agilidade, organização e controle. Um sistema informatizado e centralizado permitirá o armazenamento seguro das informações, o monitoramento em tempo real das atividades e a geração de relatórios gerenciais, beneficiando diretamente alunos, professores, orientadores e setores administrativos.

Portanto, a proposta deste trabalho é justificada pela necessidade real e urgente de modernização dos processos de estágio na FCTE, visando à melhoria da gestão acadêmica, à redução de erros e retrabalhos, e à promoção de uma experiência mais eficiente e transparente para toda a comunidade universitária.


\section{Objetivos}
Este trabalho tem como objetivo desenvolver uma aplicação web eficiente 
e centralizada para o gerenciamento de estágios na UnB. A proposta visa automatizar e simplificar processos administrativos, 
proporcionando uma plataforma que permita gestão integrada de informações, monitoramento de atividades 
e interação facilitada entre os principais atores envolvidos: estudantes, professores orientadores, empresas concedentes 
e a instituição.

Os principais \textbf{objetivos específicos} desta aplicação incluem:
\begin{itemize}
    \item Criar um sistema que permita o cadastro, acompanhamento e gestão de estágios, com interface intuitiva para os usuários;
    \item Implementar funcionalidades para controle de documentos obrigatórios, como termos de compromisso e relatórios de estágio;
    \item Assegurar um sistema robusto e seguro, que possa ser escalável e adaptado para outros departamentos ou universidades, caso necessário.
\end{itemize}